% ============================================================================
%  Chapter 3 -- Methodology
% ----------------------------------------------------------------------------
%  FORMAT CONVERSION of thesis_writing/chapter_3/chapter_3_full_v2.md
%  (frozen, scientifically approved draft; source files, in order:
%  sections_3_1_3_2_v2.md, section_3_3_v2.md, section_3_4_v1.md,
%  section_3_5_v2.md, section_3_6_v1.md, section_3_7_v1.md). Scientific
%  prose, wording, equations, notation, tables, method names and caveats are
%  preserved; only changes required for valid LaTeX have been made (math
%  delimiters, citation commands, cross-references, list/table environments,
%  labels). This is a FORMAT CONVERSION, not a rewrite -- see
%  thesis_latex/ch2_ch3_latex_integration_audit.md Part 10 for the explicit
%  content-preservation diff against the .md source.
%
%  Source editorial history (metadata, not chapter prose; relocated here
%  from the .md source's own per-section headers, exactly as Chapter 2's own
%  conversion relocated its assembly note):
%   - 3.1/3.2 (v2): grounding citation added to the opening sentence
%     (De Lange et al., 2022); "independent-per-step integration strategy" /
%     "persistent, growing-rank strategy" replaced by the names SimpleAvg /
%     RankExt, matching Chapter 2 Section~\ref{sec:lora-cl}'s naming.
%   - 3.3 (v2): two precision/editorial corrections to the SimpleAvg
%     rank-interpretation paragraph and the RankExt schedule-attribution
%     sentence (section_3_3_audit_v1.md).
%   - 3.5 (v2): the classifier weight matrix, previously bare $W$, is
%     written $W_{cls}$ throughout, to avoid colliding with Section 3.2's
%     bare $W$ (the LoRA-adapted target-layer weight, $W=W_0+\Delta W$) --
%     chapter_3_full_audit.md's notation-collision table.
%   - No fact, equation, or citation was changed by either v1->v2 pass in
%     any section.
%
%  CITATION SYSTEM: same as Chapter 2 (biblatex, style=authoryear,
%    backend=biber). All citation keys used here already exist in
%    references.bib and are already cited in Chapter 2 -- no new
%    bibliography entries were required for Chapter 3.
%
%  LABELS: semantic, ch3- prefixed, mirroring Chapter 2's sec:/fig:/tab:
%    convention. The chapter label is \label{ch:methodology} -- kept
%    identical to the placeholder already used by every existing
%    \ref{ch:methodology} in Chapter 2 (and to the ch:introduction /
%    ch:experimental-setup / ch:results / ch:discussion-conclusion pattern
%    used by every other chapter placeholder) rather than introducing a
%    differently-prefixed chap:methodology, which would require renaming
%    ~15 already-correct Chapter 2 cross-references for no functional gain.
%    See ch2_ch3_latex_integration_audit.md Part 9 for this decision.
%
%  PROCEDURES (3.7.2): rendered as enumerate lists, matching the source's
%    own numbered-step prose exactly. No algorithm/algorithmic package was
%    added -- the source is a sequential procedure description, not
%    pseudocode with loops/conditionals, so a numbered list is the least
%    intrusive faithful representation.
% ============================================================================

\chapter{Methodology}
\label{ch:methodology}

\section{Problem Formulation and Rehearsal-Free CIL Setting}
\label{sec:ch3-problem-formulation}

This thesis studies class-incremental learning as a sequence of $T$
incremental steps, $t = 1, \dots, T$ \parencite{delange2022continual}. At step $t$,
the learner is given a training set $D_t$ containing labelled examples for
a set of classes $C_t$ introduced at that step; classes are disjoint across
steps, so $C_i \cap C_j = \emptyset$ for $i \neq j$, and the cumulative
label set observed up to and including step $t$ is $C_{1:t} = C_1 \cup
\dots \cup C_t$. The partition into steps is fixed in advance and is not
revisited during training. This formulation is the class-incremental
setting already introduced in Section~\ref{sec:continual-learning}: the label space grows
monotonically, and no task identifier distinguishes which step a given
class belongs to once training has moved past it.

The constraint adopted throughout this thesis is rehearsal-free training:
no raw example from a previous step's training set $D_i$, $i<t$, is stored
or replayed at step $t$. Only the current step's data, $D_t$, is available
for gradient-based training at that step. This constraint should not be
read as a broader prohibition on carrying information forward. As
established in Section~\ref{sec:rehearsal-free}, rehearsal-free is not equivalent to a
memory-free or data-free setting; a rehearsal-free learner may still retain
model state between steps. In this thesis, the retained state can include
the frozen backbone parameters shared by all steps, adaptation parameters
produced by earlier steps, the classifier's accumulated parameters, and,
where a particular training objective calls for one, a frozen model or
classifier snapshot obtained from an earlier point in the sequence. None of
this constitutes access to raw past data: no image or label from $D_i$,
$i<t$, is ever re-examined. The specific forms this retained state takes
for each of the two compared integration strategies, and for the
distillation objective that makes use of a frozen snapshot, are given in
Sections~\ref{sec:ch3-integration-strategies} and~\ref{sec:ch3-stabilization}.

Evaluation in this setting is task-agnostic. At any point after step $t$,
the model is queried with an input belonging to one of the classes in
$C_{1:t}$ and must produce a prediction over that entire cumulative label
set, without being told which step introduced the correct class. This is
the general class-incremental inference condition and is the primary
regime in which every method compared in this thesis is assessed. A
diagnostic decomposition of this evaluation --- separating what the model
predicts when forced to compete against all classes seen so far from what
it predicts when the candidate classes are restricted to a known subset ---
is introduced later, in Section~\ref{sec:ch3-open-restricted}, once the classifier mechanism itself
has been described; it is not needed to state the basic inference
condition here.

Together, the rehearsal-free constraint and task-agnostic inference define
the methodological difficulty this thesis addresses. Each incremental step
must extend the model's competence to the new classes in $C_t$ without
degrading its competence on the classes in $C_{1:t-1}$, and it must do so
without recourse to any stored example of those earlier classes. This is
an instance of the stability--plasticity trade-off discussed in Section~\ref{sec:continual-learning},
sharpened by the absence of a replay buffer. It is this specific
difficulty --- retaining prior competence under a hard no-replay constraint,
using only retained model state --- that motivates
the two low-rank integration strategies compared in this thesis, together
with the additional stabilisation mechanisms introduced alongside them in
Section~\ref{sec:ch3-stabilization}.

To describe these mechanisms precisely, it is convenient to write the model
available after step $t$ as $f_{\theta_t}$, where the parameter set
$\theta_t$ is partitioned into three kinds of components: backbone
parameters, which are pretrained and frozen for the entire sequence and
never depend on $t$; adaptation parameters, introduced and updated by the
low-rank adaptation mechanism described in Section~\ref{sec:ch3-lora-framework}, whose organisation
across steps depends on which integration strategy is in use; and
classifier parameters, which map the shared representation onto the
cumulative label space and are updated as the label space grows. Sections~\ref{sec:ch3-integration-strategies}
and~\ref{sec:ch3-classifier-calibration} return to the adaptation and classifier components
respectively. One terminological caution is worth stating at this point,
before either strategy is introduced: because the two strategies compared
here organise their adaptation parameters differently over the incremental
sequence, a quantity as basic as ``rank'' cannot be used unqualified. The
rank configured for a single low-rank adapter, the rank added at a single
step of a persistent adapter, the cumulative rank such a persistent adapter
reaches after several steps, and the rank of a dense update obtained by
combining several independently learned adapters are, in general, four
different quantities. Section~\ref{sec:ch3-lora-framework} introduces the low-rank update itself and
fixes the first of these; the remaining distinctions are made precise once
each strategy's mechanics are described in Section~\ref{sec:ch3-integration-strategies}.

\subsection{Dataset and Protocol Instantiation}
\label{sec:ch3-protocol-instantiation}

The formulation above is deliberately independent of any particular
dataset, total class count, or number of incremental steps: it is stated
in terms of a dataset $D$, a total class count $C$, a number of steps $T$,
and a per-step class group $C_t$ of size $|C_t|$. Every method definition,
training objective, evaluation protocol, and metric given in the remainder
of this chapter is expressed in exactly these terms and does not itself
depend on any concrete choice of $C$, $T$, or $|C_t|$. This subsection
fixes the concrete instantiations used to obtain the results reported in
this thesis (Table~\ref{tab:ch3-protocol-config}); the general methodology
in the rest of this chapter is not duplicated per dataset.

The \emph{primary benchmark} is CIFAR-100, with $C=100$ classes partitioned
into five incremental steps of 20 classes each ($T=5$, $|C_t|=20$,
$5\times20$); this is the protocol used for the main results reported in
Chapter~\ref{ch:results}. Unless a passage explicitly notes otherwise,
every concrete number stated elsewhere in this chapter --- a rank schedule,
the classifier's dimensionality, and so on --- refers to this primary
protocol. A \emph{secondary protocol diagnostic} partitions the same 100
CIFAR-100 classes into four steps of 25 classes ($4\times25$,
Section~\ref{sec:ch5-ablation}). A \emph{second-dataset supporting
evaluation} uses CUB-200-2011, with $C=200$ classes in five steps of 40
classes ($5\times40$, Section~\ref{sec:ch5-cub}).

Apart from the single $4\times25$ diagnostic, no alternative protocol depth
is evaluated in this thesis; protocol depth is otherwise discussed only as
literature background (Section~\ref{sec:protocol-depth}) and as a
direction for future work (Section~\ref{sec:ch6-future-work}).

\begin{table}[htbp]
\centering
\begin{tabularx}{\textwidth}{@{}l c X@{}}
\toprule
\textbf{Dataset} & \textbf{Total classes ($C$)} & \textbf{Experimental protocol} \\
\midrule
CIFAR-100 & 100 & Primary: $5\times20$; secondary protocol diagnostic: $4\times25$ \\
CUB-200-2011 & 200 & $5\times40$ (second-dataset supporting evaluation) \\
\bottomrule
\end{tabularx}
\caption[Dataset and protocol configuration]{Datasets and class-incremental
protocols used in this thesis. No other part of this chapter's methodology
depends on the contents of this table.}
\label{tab:ch3-protocol-config}
\end{table}

\section{LoRA Adaptation Framework}
\label{sec:ch3-lora-framework}

Both integration strategies compared in this thesis adapt the same frozen
pretrained backbone. At the level of backbone LoRA adaptation, the two
strategies share the same low-rank update mechanism and differ in how those
parameters are organised and integrated across steps. This section defines the shared
mechanism; the two strategies themselves are introduced in Section~\ref{sec:ch3-integration-strategies}.

The backbone is the vision encoder of CLIP ViT-B/16 \parencite{radford2021learning},
a Vision Transformer \parencite{dosovitskiy2021image}. Its parameters are frozen
throughout training, at every step and for both compared strategies: no
gradient updates the pretrained weights themselves. The pooled
representation it produces is passed to a linear classifier, described
separately in Section~\ref{sec:ch3-classifier-calibration}. Adaptation is confined entirely to a small set
of additional low-rank parameters injected into the backbone, following
the low-rank adaptation (LoRA) formulation of \cite{hu2021lora}.

For a frozen weight matrix $W_0 \in \mathbb{R}^{d_{out} \times d_{in}}$ of
a target linear layer, LoRA leaves $W_0$ unchanged and represents the
adaptation as an additive low-rank term:

\begin{equation}
\label{eq:ch3-lora-update}
W = W_0 + \Delta W, \qquad \Delta W = s \, B A,
\end{equation}

where $A \in \mathbb{R}^{r \times d_{in}}$ and $B \in \mathbb{R}^{d_{out}
\times r}$ are the trainable low-rank factors, $r$ is the rank of the
update, and $s$ is a scaling factor applied to the product. Because
$\Delta W$ is formed as the product of two rank-$r$ factors, a single such
update has rank at most its configured rank $r$. The scaling factor is set as the ratio of a
configured constant $\alpha$ to the rank, $s = \alpha / r$, so that the
effective contribution of the adapter can be held fixed even if $r$ is
later changed. In the implementation studied here, LoRA is applied to the
query and value projection matrices of every attention block in the
backbone; this is the sole set of weights either compared strategy is
permitted to adapt, and it is the same set for both --- a point worth stating
explicitly, since it is the kind of configuration detail under which the
two strategies could otherwise be made to differ without that difference
being part of either method's definition.

Within this shared framework, SimpleAvg (Section~\ref{sec:ch3-simpleavg}) is configured with
rank $r = 80$ and $\alpha = 160$, giving a scaling of $s = \alpha/r = 2$.
RankExt (Section~\ref{sec:ch3-rankext}) organises its low-rank parameters differently
across the incremental sequence --- a distinction deferred in full to that
section --- but is configured so that its effective scaling is the same
constant, $s = 2$, at every step regardless of how much rank it has
accumulated. Using the same effective LoRA scaling for both strategies
ensures that observed differences between them are not attributable simply
to a difference in LoRA scaling.

Rank itself, however, is not a single quantity once adaptation is spread
across several incremental steps, and it is worth fixing the vocabulary
before either strategy is described. The rank $r$ fixed above for a single
LoRA factorisation is its \emph{configured} rank. A persistent adapter that
accumulates rank over the sequence has, at a given step, an \emph{incremental}
rank --- the amount added at that step --- and a \emph{cumulative} rank --- the
running total after that step; Section~\ref{sec:ch3-rankext} gives both a precise
definition. A separate quantity again is the rank of a \emph{dense combination}
of several independently learned low-rank updates: summing or averaging
several rank-$r$ updates does not, in general, produce a matrix whose rank
is still bounded by $r$. Concretely, and stated here only as a bound on
notation rather than as a result: each SimpleAvg adapter has configured
rank 80, but when several such independently learned dense updates are
later combined into one, the resulting combined update is not thereby
constrained to rank 80 as well.
The exact combination rule, and what this implies for the comparison
between the two strategies, is developed in Section~\ref{sec:ch3-integration-strategies} and revisited,
with supporting evidence, in later chapters.

The remainder of this chapter builds on this shared framework. Section~\ref{sec:ch3-integration-strategies}
introduces the two ways in which this thesis organises low-rank adaptation
across a continual sequence --- an independent adapter retrained at every
step and combined afterwards, denoted SimpleAvg, and a persistent adapter
whose rank is extended step by step, denoted RankExt --- and shows precisely
how each retains state consistently with the rehearsal-free constraint of
Section~\ref{sec:ch3-problem-formulation}.

\section{Continual LoRA Integration Strategies}
\label{sec:ch3-integration-strategies}

Both strategies compared in this thesis adapt the frozen backbone using
exactly the low-rank update defined in Section~\ref{sec:ch3-lora-framework}; nothing in the
mechanism itself differs between them. What differs is how the low-rank
parameters introduced at successive steps are retained and integrated over
the continual sequence. SimpleAvg trains an independent adapter at each
step and combines the resulting updates by dense averaging once training
is complete. RankExt instead maintains a single persistent adapter whose
rank grows over the sequence: the low-rank blocks learned at earlier steps
are kept and frozen, while a new block is appended and trained at each
step. The two subsections below define each strategy in turn; no claim
about their relative empirical performance is made here.

\subsection{SimpleAvg}
\label{sec:ch3-simpleavg}

At each step $t$, SimpleAvg trains a fresh LoRA adapter on a fresh copy of
the frozen pretrained backbone, entirely independently of the adapters
trained at any other step. Independence here is architectural, not only a
matter of not sharing gradients: step $t$'s adapter is attached to the
original pretrained weights, not to a backbone already carrying the
adaptation learned at step $t-1$, so no step begins from an accumulated
or partially adapted state. No SimpleAvg adapter parameters are carried
forward as trainable state from one step to the next. Consistently with
Section~\ref{sec:ch3-problem-formulation}'s rehearsal-free
constraint, no raw example from an earlier step's training set is
available while step $t$'s adapter is being trained. The adapter is
configured as fixed in Section~\ref{sec:ch3-lora-framework}: rank $r=80$, $\alpha=160$, giving
scaling $s=\alpha/r=2$, applied to the query and value projection matrices
of every attention block. Each step also trains its own classifier head
from a fresh initialisation, over the full label space but with only that
step's classes ever positive; how this head is later reconciled with the
heads trained at other steps is addressed below.

For each adapted layer $m$, training step $t$ produces low-rank factors
$A_t^{(m)}, B_t^{(m)}$, from which the per-step dense update
$\Delta W_t^{(m)}$ is reconstructed using the LoRA rule of
Equation~\eqref{eq:ch3-lora-update}. It is this dense update, not the factors $A_t^{(m)}$ and $B_t^{(m)}$
themselves, that SimpleAvg later combines across steps; the factors are
never averaged directly. Once all $T$ steps have been trained, the final
model is formed by taking the elementwise arithmetic mean of the $T$
per-step dense updates for each adapted layer,

\begin{equation}
\label{eq:ch3-simpleavg-merge}
\Delta \bar{W}^{(m)} = \frac{1}{T} \sum_{t=1}^{T} \Delta W_t^{(m)},
\qquad
W_{final}^{(m)} = W_0^{(m)} + \Delta \bar{W}^{(m)},
\end{equation}

and this mean --- not a sum --- is added to the frozen base weight to obtain
the final adapted model. This combination is performed exactly once, after
the entire continual sequence has been trained; SimpleAvg does not rebuild
a merged model after every step, and no intermediate merged model is ever
evaluated as the method's output. A separate, training-time-only
artifact --- a frozen snapshot obtained by merging only the steps already
trained so far, rebuilt at the start of each subsequent step --- is used as
a teacher for one of the stabilisation mechanisms introduced in Section~\ref{sec:ch3-kd}.
This snapshot is distinct from, and computed separately from, the
single final evaluation merge described above, and plays no further role
in this section.

The classifier that accompanies the final merged backbone is not itself
averaged. Each step trains its own head, as noted above, but only the rows
belonging to that step's classes are ever meaningfully trained; the final
classifier is assembled by copying, for every class, the weight row and
bias from the step that introduced it, so that classifier assembly is, by
construction, a separate operation from the LoRA-delta averaging described
above. The full classifier mechanism, including how these rows interact
with classifier rescaling, is developed in Section~\ref{sec:ch3-classifier-calibration}.

A configured rank of 80 describes a single SimpleAvg adapter, not the
final combined update: the final dense merged update is not constrained to
the configured rank of a single rank-80 adapter. This combination rule is a direct instantiation,
for low-rank
adapters, of the unweighted dense-averaging baseline already positioned
against the broader model-merging literature in Section~\ref{sec:model-merging}: SimpleAvg is
presented here as a deliberate baseline, independent per-step training
followed by post-hoc dense averaging with no clustering, weighting, or
orthogonalisation of the contributions being combined, included to give a
reference point for the persistent, growing-rank strategy described next.
It is not put forward as a methodological contribution of this thesis.

\subsection{Rank Extension (RankExt)}
\label{sec:ch3-rankext}

RankExt organises the same low-rank adaptation mechanism differently.
Rather than training a fresh adapter at every step, it maintains a
persistent adapter state across the sequence. At each step, the model is
reconstructed from the pretrained backbone and the cumulative adapter state
carried forward from the previous step; previously learned blocks are
reinstated as frozen, active components, and one new trainable rank block
is appended. The pretrained base weight $W_0$ remains frozen throughout, as
it does for SimpleAvg.

Write the cumulative rank reached after step $t$ as $R_t$ (with $R_0=0$),
and the block of factors learned at step $k$ as $A_k, B_k$, of incremental
rank $R_k - R_{k-1}$. The cumulative update at step $t$ is then the sum of
the contributions of every block learned so far,

\begin{equation}
\label{eq:ch3-rankext-cumulative}
\Delta W_t = s \sum_{k=1}^{t} B_k A_k,
\end{equation}

where blocks $k<t$ are frozen --- they receive no gradient and are not
included in any optimiser's parameter group --- and only block $k=t$ is
trainable, with the same constant scaling $s=2$ fixed in Section~\ref{sec:ch3-lora-framework}
applied uniformly to every block regardless of when it was learned. The
implementation obtains this sum by concatenating the frozen and new factor
matrices rather than by literally summing block-wise products, but the two
are equivalent for a linear map: writing $A^{(t)}$ as the row-concatenation
of all blocks up to $t$ and $B^{(t)}$ as their column-concatenation,
$B^{(t)}A^{(t)} = \sum_{k=1}^{t} B_k A_k$.

In the primary $5\times20$ protocol (Section~\ref{sec:ch3-protocol-instantiation})
used for the main comparison, the cumulative RankExt schedule is
$R_t=[16,32,48,64,80]$ for $t=1,\dots,5$, corresponding to an increment of
rank 16 at each step. This is the schedule against which RankExt is
compared to SimpleAvg in the main results of this thesis; it is a
configurable property of the architecture, not a consequence of the
general formulation above. A secondary $4\times25$ protocol diagnostic
uses an adjusted schedule $[20,40,60,80]$ so that the final cumulative rank
remains 80; that experiment is described in Chapter~\ref{ch:results} and is
not treated as a controlled single-factor protocol-depth study. A previously learned block's parameters have zero gradient from
the step after it was learned onward, contribute to the forward pass
exactly as before, and are verified to remain unchanged for the remainder
of training; freezing is total, not partial. A newly appended block's $A$
factor is initialised as any fresh LoRA factor would be, while its $B$
factor is initialised to zero, so that a block contributes nothing to the
forward pass before it has been trained.

RankExt's effective scaling remains the constant $s=2$ at every step,
independently of how much cumulative rank has been reached. The
implementation derives, purely for internal bookkeeping, an $\alpha$-like
quantity proportional to the cumulative rank so that this ratio is
preserved as the rank grows; this derived quantity is not an independently
tuned hyperparameter and should not be read as one.

The two architectures are therefore configured with an explicit and
intentional asymmetry in how much adaptation capacity is introduced at
each step: SimpleAvg trains a fresh rank-80 adapter at every step, while
the canonical RankExt configuration --- in the primary $5\times20$ protocol
--- adds rank 16 per step and reaches cumulative rank 80 only at the final
step. This asymmetry is disclosed
here as part of the architectural definition; its experimental
consequences are not interpreted in this chapter, and additional
capacity-sensitivity configurations of RankExt are examined later in this
thesis. Because RankExt's cumulative update at step $t$ is formed from
concatenated factors of total rank $R_t$, its rank cannot exceed $R_t$,
and hence, under the canonical schedule, cannot exceed 80.

As with SimpleAvg, the classifier is not folded into the low-rank update:
a single classifier persists across the whole sequence, with rows for
previously introduced classes carried forward unchanged and only the rows
for the current step's classes updated at that step. The full row-training
and classifier-rescaling mechanism is again deferred to Section~\ref{sec:ch3-classifier-calibration}. Growing-rank,
persistent low-rank adaptation of this kind is not introduced by this
thesis; it instantiates a design already present in the continual-LoRA
literature discussed in Section~\ref{sec:lora-cl}. What this thesis contributes is the
controlled comparison between this persistent strategy and the
independent-averaging strategy of Section~\ref{sec:ch3-simpleavg}, together with the
stabilisation mechanisms introduced next, in Section~\ref{sec:ch3-stabilization}.

To summarise the contrast that runs through the remainder of this thesis:
SimpleAvg learns independently at each step, reconstructs a dense update
per step, and averages these updates once after the full sequence has been
trained; RankExt instead retains every previously learned block, freezes
it permanently, appends one new trainable block per step, and carries
forward a single evolving adapter state throughout the sequence. Both strategies are, in
Section~\ref{sec:ch3-stabilization}, paired with additional stabilisation mechanisms that act on
top of this base architecture.

\section{Continual Stabilisation Mechanisms}
\label{sec:ch3-stabilization}

Neither integration strategy defined in Section~\ref{sec:ch3-integration-strategies} places any explicit
constraint on how much the shared representation is allowed to change when
a new group of classes is learned: SimpleAvg's independence between steps
and RankExt's frozen prior blocks limit \emph{where} new parameters can act, but
not how far training is free to push the current step's own contribution.
This thesis studies two additional mechanisms layered on top of either
architecture --- knowledge distillation and FactorOrth regularisation,
whose concrete formulation is family-specific --- and their combination. The same stabilisation
principles are examined within both architectural families of
Section~\ref{sec:ch3-integration-strategies}, with family-specific
implementations where required by the architecture.
Neither mechanism is presented here as
universally necessary or as empirically superior to the other; their
measured effect, individually and combined, is examined in Chapter~\ref{ch:results}.

\subsection{Knowledge Distillation}
\label{sec:ch3-kd}

Knowledge distillation trains the current step's model to match the
temperature-softened output distribution of a frozen teacher, rather than
only the hard label, so that the relative scores the teacher assigns to
non-target classes also shape the student's training signal \parencite{hinton2015distilling}.
This is a form of self-distillation between successive states of the same
model family, aimed at stabilising continual learning rather than
compressing a larger teacher into a smaller student. In this
implementation, distillation is computed only on the current step's own
input batch --- no image or label from an earlier step is replayed for this
purpose, consistent with the rehearsal-free constraint of Section~\ref{sec:ch3-problem-formulation} ---
and its scope is family-specific. For RankExt, it is applied to the
\textbf{full} $C$-way output distribution, which covers every class seen
so far, including the classes newly introduced at the current step. For SimpleAvg, it is
restricted to the previously seen classes: teacher and student logits are
both restricted to the classes introduced at earlier steps before the
softened distributions are formed. Writing $z_t$ and $z_s$ for the
teacher's and student's logits on the same batch (restricted to the
previously seen classes in the SimpleAvg case), and $\tau$ for the
distillation temperature,

\begin{equation}
\label{eq:ch3-kd}
p_t = \operatorname{softmax}(z_t/\tau), \qquad
p_s = \operatorname{softmax}(z_s/\tau), \qquad
\mathcal{L}_{KD} = \tau^2 \, D_{KL}(p_t \,\|\, p_s),
\end{equation}

with the teacher's distribution $p_t$ computed without gradient and used as
the fixed reference --- the student is drawn toward the teacher, never the
reverse. The temperature is $\tau=2$, and the canonical distillation weight
is $\kappa=1.0$. Because a teacher requires at least one already-trained
step, distillation is active only from the second step onward. For
RankExt, its weight and temperature remain constant whenever it is active;
for SimpleAvg, the distillation weight is ramped linearly over the first
epoch of each step, as specified in Chapter~\ref{ch:experimental-setup}.

The two integration strategies differ in what the teacher is, for
architectural reasons rather than by an independent design choice. Because
SimpleAvg has no single evolving model until the sequence is complete, its
teacher cannot simply be ``the model as it stood after the previous step'' ---
no such object exists during training. Instead, the teacher is a frozen
snapshot produced by applying, to only the steps already trained, exactly
the dense-averaging merge and per-class classifier row assembly defined in
Section~\ref{sec:ch3-simpleavg} --- the same combination rule, evaluated over a growing prefix
of the sequence rather than the complete one, and rebuilt from scratch at
the start of each subsequent step because it is not itself stored between
steps. This training-time snapshot is distinct from the single merge
performed once after the complete sequence, described in Section~\ref{sec:ch3-simpleavg}. RankExt, by contrast, carries a complete adapter and classifier state
forward at every point in the sequence, so its teacher is simply the model
defined by that state at the end of the previous step, before the current
step's new block is added: no separate merge-and-stitch operation is
needed to construct it, because RankExt's own state already defines a
directly usable model at every step boundary.

This mechanism should not be described as replay: no raw example from an
earlier step is stored or reused, only the current step's own images are
passed through the frozen teacher. Nor should it be described as
data-free --- a live, if frozen, teacher model is required at every
distillation step, and the constraint this thesis satisfies concerns which
data that teacher is not given, not the absence of data altogether.
Distillation's temperature is fixed throughout; the orthogonality
regularisation this thesis studies alongside it is defined next.

\subsection{FactorOrth Regularisation}
\label{sec:ch3-factororth}

FactorOrth denotes the orthogonality-based regularisation used in both
integration families, but its concrete formulation is family-specific.
Each family's formulation follows from what that family retains from
earlier steps: SimpleAvg retains a reconstructed dense update per step,
whereas RankExt retains its earlier low-rank factors as frozen blocks. The
two formulations are therefore two different regularisation operators that
share a common label and a common purpose --- discouraging the current
step's update from coinciding with directions already used by earlier
steps --- and neither should be described in terms of the other. In both
families the penalty is inactive at the first step, for which no earlier
update exists, and its coefficient is ramped in linearly over the first
epoch of each step, as specified in Chapter~\ref{ch:experimental-setup}.

\textbf{A. SimpleAvg FactorOrth (dense-update formulation).} Because each
SimpleAvg step already reconstructs its dense update and retains it for the
final merge, the SimpleAvg formulation compares the current step's dense
update $\Delta W^{(m)}$ directly with each earlier step's retained dense
update $\Delta W_u^{(m)}$, rather than comparing low-rank factors:
\begin{equation}
\label{eq:ch3-factororth-simpleavg}
\mathcal{L}_{orth,\,SA}^{(m)} = \frac{1}{t-1}\sum_{u<t}
\left(\frac{\langle \Delta W^{(m)}, \Delta W_u^{(m)}\rangle_F}
{\lVert \Delta W^{(m)}\rVert_F \,\lVert \Delta W_u^{(m)}\rVert_F}\right)^{2},
\end{equation}
where $\langle\cdot,\cdot\rangle_F$ is the Frobenius inner product. The
SimpleAvg FactorOrth loss is this quantity averaged over the adapted
layers. Each earlier step's dense update enters individually, not through
an averaged reference, and is held fixed, so only the current step's
factors receive gradient. In the run behind the results of
Chapter~\ref{ch:results}, its regularisation strength is $\lambda_{orth}=20$.

\textbf{B. RankExt FactorOrth (factor-space formulation).} The RankExt
formulation regularises the geometry of the low-rank factors learned at the
current step relative to the frozen, accumulated prior rank blocks
$A_{ref}, B_{ref} = A_{frozen}^{(t)}, B_{frozen}^{(t)}$ already introduced
in Section~\ref{sec:ch3-rankext}, computed directly from normalised factors
rather than from the reconstructed weight update. For a target layer, with
the new block's factors $A, B$, each row of $A$ and $A_{ref}$ is normalised
to unit length, and each column of $B$ and $B_{ref}$ is normalised to unit
length:

\begin{equation}
\label{eq:ch3-factororth-a}
\hat A_{ref} = \operatorname{rownorm}(A_{ref}), \quad
\hat A = \operatorname{rownorm}(A), \quad
O_A = \hat A_{ref}\hat A^{\top}, \quad
\phi_A = \lVert O_A \rVert_F^2,
\end{equation}
\begin{equation}
\label{eq:ch3-factororth-b}
\hat B_{ref} = \operatorname{colnorm}(B_{ref}), \quad
\hat B = \operatorname{colnorm}(B), \quad
O_B = \hat B_{ref}^{\top}\hat B, \quad
\phi_B = \lVert O_B \rVert_F^2,
\end{equation}

so that $O_A$ and $O_B$ are cross-factor overlap matrices between the
reference and current directions, and $\phi_A,\phi_B$ are the squared
Frobenius norms of those overlaps. The per-layer penalty is
$\phi_A+\phi_B$, and the RankExt FactorOrth loss is this quantity averaged
over the layers the penalty is computed at. Because every direction
entering $O_A$ and $O_B$ has first been normalised to unit length, the
penalty discourages the new block's directions from \emph{coinciding} with
the frozen reference directions, regardless of either side's original
magnitude. The reference is not an average: it is the architecture's own
already-frozen state, used directly as a fixed target. Its regularisation
strength is $\lambda_{orth}=50$; the coefficient is therefore not shared
across families.

The role of the orthogonality term also differs between the two
strategies, as a consequence of the architectural difference established in
Section~\ref{sec:ch3-integration-strategies}. RankExt's prior blocks are
already hard-frozen by construction, so FactorOrth there is not protecting
those blocks from being overwritten --- it instead discourages the newly
appended block from learning directions that merely duplicate ones the
frozen blocks already occupy. SimpleAvg has no such architectural freezing,
since every step's adapter is trained independently; there, FactorOrth is
the only mechanism discouraging the current step's dense update from
aligning with those of earlier steps before the dense updates are ever
combined.

FactorOrth does not originate the idea of an orthogonality penalty between
old and new low-rank updates; that idea, and this thesis's positioning of
its own formulations relative to it --- including the related structural
alternative to a soft penalty --- are established in Section~\ref{sec:orthogonality}.
What is specific to this implementation is the exact form of the two
family-specific formulations and the references they use.

Every variant is trained with a standard cross-entropy loss on the current
step's labelled batch as its base term, to which the mechanisms of
Sections~\ref{sec:ch3-kd} and~\ref{sec:ch3-factororth} add further terms
whenever they are active. The exact variant-specific loss composition,
auxiliary stabilisation terms, and warmup schedules used in the experiments
are specified in Chapter~\ref{ch:experimental-setup}.

\section{Classifier Construction and Rescaling}
\label{sec:ch3-classifier-calibration}

Representation integration, as defined in Section~\ref{sec:ch3-integration-strategies}, and the classifier
that turns that representation into class predictions are handled as two
separate operations on two separate sets of parameters. Nothing in
SimpleAvg's dense-averaging merge or RankExt's block-sum forward pass
touches the classifier, and, symmetrically, the classifier rescaling
described below never touches the backbone or the low-rank adaptation
parameters. This section defines both operations in turn.

\subsection{Classifier Construction}
\label{sec:ch3-classifier-construction}

The classifier is a single linear layer that maps the shared pooled
representation $h$ to logits $z$ through a weight matrix
$W_{cls}\in\mathbb{R}^{C\times d}$ and a bias vector $b\in\mathbb{R}^{C}$,
covering every class the model will ever see ($C=100$ for CIFAR-100 and
$C=200$ for CUB-200-2011, Section~\ref{sec:ch3-protocol-instantiation}). This is precisely the classifier component of the three-way
parameter partition of $f_{\theta_t}$ introduced in Section~\ref{sec:ch3-problem-formulation}, alongside
the frozen backbone and the adaptation parameters of Section~\ref{sec:ch3-integration-strategies}. All $C$
rows of $W_{cls}$ and $b$ exist from the
first step onward, for both strategies; what changes across steps and
across strategies is which rows are meaningfully trained at a given point,
not the tensor's shape, and in both cases the set of rows available to
compete at inference tracks exactly the cumulative label set $C_{1:t}$
already defined in Section~\ref{sec:ch3-problem-formulation}. No row is normalised during training --- the
normalisation applied in Section~\ref{sec:ch3-calibration} is strictly post-hoc.

For SimpleAvg, each step trains its own classifier alongside that step's
adapter, as already noted in Section~\ref{sec:ch3-simpleavg}. The final classifier used with
the single averaged LoRA update is not obtained by averaging these
per-step classifiers. Instead, for every class, the final weight row and
its corresponding bias entry are both copied from whichever step
introduced and trained that class: the final classifier is assembled by
copying each class row and bias from the step that trained that class,
not by averaging the per-step classifiers. This assembly operation is
independent of the
dense-delta averaging described in Section~\ref{sec:ch3-simpleavg}: the two are combined into one
final model, but neither is derived from the other.

For RankExt, the classifier state is carried forward across the entire
sequence, as already noted in Section~\ref{sec:ch3-rankext}, mirroring the
persistence of the low-rank blocks themselves: when the model is
reconstructed at each step, both the classifier weights and the adapter
factors learned so far are reinstated rather than reinitialised. Rows belonging to classes
introduced at earlier steps are carried forward unchanged into every later
step, and only the rows for the current step's classes receive gradient at
that step. The implementation enforces this guarantee redundantly ---
through a gradient mask during training and a hard restoration of
protected rows afterward --- but the property that matters methodologically
is simply stated: old class rows are preserved exactly, and only new class
rows are trained.

\subsection{Classifier Rescaling}
\label{sec:ch3-calibration}

After training is complete, and before either evaluation mode introduced
in Section~\ref{sec:ch3-open-restricted} is computed, the final classifier's weight rows are
rescaled by a post-hoc procedure applied identically to both strategies.
Classifier rescaling is not part of gradient-based training, does not touch the
backbone or the low-rank adaptation parameters, and never modifies the
bias vector $b$ --- only the weight rows $W_{cls}[c,:]$ are rescaled, based on
their own $\ell_2$ norm, $n_c=\lVert W_{cls}[c,:]\rVert_2$. The rescaling is
norm-based: each step's block of rows is multiplied by a single positive
factor derived from a target row norm for that step, so that an entire
step's rows are scaled together rather than individually.

This family of rescaling rules is conceptually related to the norm-balancing
principle of Weight Aligning \parencite{zhao2020maintaining}, which observes
that later-trained class rows tend to have larger norms than earlier ones
and corrects for it by rescaling the new-class rows toward the old classes'
mean norm. The classifier rescaling used in this thesis is not a direct
implementation of Weight Aligning: several target-norm variants are
implemented, differing in how the target norm for a step is chosen and
whether it is adjusted by that step's own training convergence; their exact
form, and which variant is used for the results reported in this thesis,
are configuration detail specified in Chapter~\ref{ch:experimental-setup}
rather than part of the methodological definition given here.

Because every variant rescales an entire step's block of rows by one
common positive factor, and does so identically regardless of which
evaluation mode will later be applied, the rescaled classifier is what
both open and restricted evaluation are computed from. Section~\ref{sec:ch3-open-restricted}
returns to this property once both evaluation modes have been defined.

\section{Open and Restricted Evaluation Protocols}
\label{sec:ch3-open-restricted}

For any fixed checkpoint and test example, open and restricted evaluation
use the same model, forward pass, and logits; only the set of logits
eligible for the final argmax differs. The paired open/restricted headline
diagnostic reported for the completed model uses the final model after classifier rescaling
(Section~\ref{sec:ch3-classifier-calibration}). Trajectory-based
measurements --- backward transfer and forgetting
(Section~\ref{sec:ch3-bwt-forgetting}) --- additionally require
step-specific references: intermediate RankExt checkpoints, and
SimpleAvg's pre-merge specialists. Two different views are taken of a
checkpoint's logits. The first, open evaluation, lets a prediction
compete against every class the model has learned. The second, restricted
evaluation, narrows the competition to only the classes introduced
alongside the example being scored. Nothing about the model, the
classifier, the examples, or the logits themselves differs between the two
views; only which logits are eligible for the final argmax changes. This
section defines both precisely, since the distinction between them bears
directly on how the results in later chapters should be read.

Open evaluation is this thesis's term for what the continual-learning
literature calls task-agnostic evaluation \parencite{masana2020classincremental}: the
predicted class for an input $x$ is the argmax of the classifier's logits
$z(x)$ over the complete cumulative label set observed by the point at
which evaluation is performed, with no task identifier supplied,

\begin{equation}
\label{eq:ch3-open-eval}
\hat y_{open}(x) = \arg\max_{c\,\in\,C_{1:t}} z_c(x).
\end{equation}

In this implementation, every open accuracy reported for the completed
model --- whether computed over the classes of a single step or over every
class the model has been trained on --- is obtained from the single final
model produced after the complete sequence of $T$ steps, at which point
every class this thesis studies has already been introduced. For these
final results the open argmax is therefore taken with $t=T$, i.e.\ over
$C_{1:T}$, the complete set of classes seen across the sequence; the
step-specific accuracies used only for backward transfer and forgetting
are computed from the references defined in
Section~\ref{sec:ch3-bwt-forgetting}.

Restricted evaluation is this thesis's term for task-oracle evaluation
\parencite{masana2020classincremental}. For an example belonging to the group of classes
introduced at step $t$, $C_t$, the same logits $z(x)$ are used, but every
logit for a class outside $C_t$ is excluded from the argmax:

\begin{equation}
\label{eq:ch3-restricted-eval}
\hat y_{restr}(x) = \arg\max_{c\,\in\,C_t} z_c(x).
\end{equation}

Task identity enters only here, and only to determine which classes
constitute $C_t$ for the example being scored --- it is not used for open
evaluation, and even for restricted evaluation it never selects a
different model, a different classifier, or a different forward pass. The
model is not retrained, the classifier is not replaced, and the
representation is not modified; the same logits computed once are simply
read under a narrower mask.

Because open evaluation forces a prediction to compete against every class
the model has learned, it reflects two things at once: how well the model
discriminates within the group of classes an example actually belongs to,
and how well that group's classes compete against every other group under
a single argmax. Restricted evaluation removes cross-group competition by
construction, providing a more direct view of within-group discrimination.
Comparing the two accuracies for the
same group of examples is therefore informative about, and can help
diagnose, whether a given pattern in the results is more consistent with a
limitation in within-group discrimination or with cross-group competition
at the output layer; it does not, on its own, prove which of the two is
responsible in any particular case. Open evaluation remains the primary
regime in which this thesis's methods are assessed, since it is the one
that matches the task-agnostic inference condition set out in Section~\ref{sec:ch3-problem-formulation};
restricted evaluation is a supplementary diagnostic, not an alternative
headline metric, and should never be read as recovering a model's ``true''
or ``actual'' retention --- it reports accuracy once one specific source of
competition has been removed, nothing more.

This diagnostic role connects directly to the classifier rescaling of
Section~\ref{sec:ch3-calibration}. Because the same positive factor is
applied to all weight rows within a step, classifier rescaling does not
alter their relative weight-vector scaling. However, classifier biases are
left unchanged, so restricted accuracy is not mathematically guaranteed to
remain invariant under the rescaling. Open accuracy is more directly
exposed to classifier rescaling, since its argmax compares logits across
groups that the rescaling may have scaled differently.

\section{Method Variants and Overall Continual-Learning Procedure}
\label{sec:ch3-method-variants}

The components defined in Sections~\ref{sec:ch3-integration-strategies}--\ref{sec:ch3-open-restricted} --- the two integration
strategies, knowledge distillation, orthogonality regularisation, classifier
construction and rescaling, and the open/restricted evaluation
protocols --- are instantiated as a controlled set of eight method variants.
Their purpose is to separate the contribution of the continual integration
strategy from that of the two stabilisation mechanisms while holding the
rest of the pipeline --- the backbone, the classifier mechanism, the
classifier rescaling, and the evaluation protocols --- fixed across all
eight. This section assembles those variants, describes the end-to-end
procedure each family follows, and defines the two continual-learning
metrics this thesis reports; it does not compare their outcomes.

\subsection{Eight Method Variants}
\label{sec:ch3-eight-variants}

Within each integration strategy, knowledge distillation and FactorOrth
regularisation are switched on or off, giving four variants per family and
eight in total. Every variant shares the same frozen backbone, the same
target modules, the same classifier mechanism, the same classifier
rescaling, and the same evaluation protocols (Sections~\ref{sec:ch3-lora-framework}, \ref{sec:ch3-classifier-calibration}, \ref{sec:ch3-open-restricted}).
The variants are organised conceptually as a two-by-two KD/FactorOrth
grid within each integration family, but variant-specific warmups and
RankExt-specific auxiliary protection mechanisms mean that the complete
experiment is not a strictly factorial design. ``FactorOrth'' is a common
label: its implementation, and its coefficient, differ by family
(Section~\ref{sec:ch3-factororth}).
Table~\ref{tab:ch3-eight-methods} sets out the resulting conceptual grid.

\begin{table}[htbp]
\centering
\footnotesize
\begin{tabularx}{\textwidth}{@{}>{\raggedright\arraybackslash}p{4cm} >{\raggedright\arraybackslash}p{4.5cm} c c@{}}
\toprule
\textbf{Method} & \textbf{Integration strategy} & \textbf{KD} & \textbf{FactorOrth} \\
\midrule
SimpleAvg & independent-per-step, dense mean-merge (Section~\ref{sec:ch3-simpleavg}) & -- & -- \\
SimpleAvg + KD & same & \checkmark & -- \\
SimpleAvg + FactorOrth & same & -- & \checkmark \\
SimpleAvg + FactorOrth + KD & same & \checkmark & \checkmark \\
RankExt & persistent growing-rank (Section~\ref{sec:ch3-rankext}) & -- & -- \\
RankExt + KD + Protect & same & \checkmark & -- \\
RankExt + FactorOrth & same & -- & \checkmark \\
RankExt + FactorOrth + KD + Protect & same & \checkmark & \checkmark \\
\bottomrule
\end{tabularx}
\caption[Eight method variants]{The eight method variants compared in this
thesis, organised as a conceptual grid of knowledge distillation (KD) and
FactorOrth regularisation within each of the two integration strategies of
Section~\ref{sec:ch3-integration-strategies}. FactorOrth is a common label
whose formulation is family-specific: dense-update orthogonality for
SimpleAvg and factor-space orthogonality for RankExt
(Section~\ref{sec:ch3-factororth}). KD-bearing RankExt variants are
additionally coupled to projected feature protection
(Section~\ref{sec:ch4-rankext-auxiliary}), hence ``+ Protect''.}
\label{tab:ch3-eight-methods}
\end{table}

Within RankExt, activating knowledge distillation does not only add the
distillation loss term of Section~\ref{sec:ch3-kd}: the implementation also
couples it to further, RankExt-specific configuration choices, specified in
Chapter~\ref{ch:experimental-setup}, that are not independent of whether
distillation is active. In particular, the implemented KD-bearing RankExt
configurations are coupled to projected feature protection, as specified
in Section~\ref{sec:ch4-rankext-auxiliary}; later result labels therefore
include ``+ Protect'' (for example, RankExt + FactorOrth + KD + Protect) to
make this implementation detail visible. Protect is not an independent
toggle in the method set. SimpleAvg carries no such coupling, since it has
no persistent block for those choices to act on. The eight-method set as a
whole is therefore best described as a structured family-wise ablation over
integration strategy, KD, and orthogonality regularisation, rather than as
a strictly factorial experiment.

\subsection{Overall Continual-Learning Procedure}
\label{sec:ch3-procedure}

Both families share the same outer loop over $T$ incremental steps and the
same three closing operations, but differ in what is retained between
steps and in how the final model is assembled. Each is described
separately below, followed by the shared closing sequence.

\textbf{Procedure (SimpleAvg).} For each step $t=1,\dots,T$:

\begin{enumerate}
\item Receive the current step's data $D_t$ and its class group $C_t$
   (Section~\ref{sec:ch3-problem-formulation}); no raw example from any earlier $D_i$, $i<t$, is
   available.
\item Configure the current adaptation state: start a fresh copy of the
   frozen pretrained backbone (Section~\ref{sec:ch3-lora-framework}) and inject a new LoRA adapter
   and a new classifier head --- never continuing from a previous step's
   adapted weights (Section~\ref{sec:ch3-simpleavg}).
\item Train on $D_t$ with the classification objective and whichever
   stabilisation mechanisms are active for the selected variant
   (Section~\ref{sec:ch3-stabilization}) --- the distillation term against a
   frozen running merge of the steps already trained, and/or the
   dense-update FactorOrth term against each earlier step's retained dense
   update.
\item Update the classifier state by training the current step's head, as
   part of the same training loop as step 3 (Section~\ref{sec:ch3-classifier-construction}).
\item At the end of the step, retain the dense update and trained classifier
   rows required for final assembly; the same retained per-step dense
   updates also supply the running merge used as the distillation teacher
   and, for FactorOrth-bearing variants, the reference against which the
   next step's regularisation is computed. The step's own model instance,
   backbone copy, and adapter object are not themselves carried forward.
   No raw training example is retained.
\end{enumerate}

After all $T$ steps, the final SimpleAvg model used for evaluation is
assembled by averaging the retained dense updates and stitching the
classifier rows (Sections~\ref{sec:ch3-simpleavg}/\ref{sec:ch3-classifier-construction}).
Temporary running merged snapshots may additionally be constructed during
training when required by a stabilisation mechanism such as knowledge
distillation (Section~\ref{sec:ch3-kd}); the object described here is the
single final model produced once, after training, and used for evaluation.

\textbf{Procedure (RankExt).} For each step $t=1,\dots,T$:

\begin{enumerate}
\item Receive $D_t$ and $C_t$; as for SimpleAvg, no raw example from any
   earlier step is available.
\item Reconstruct the current model from the pretrained backbone and the
   persistent adapter and classifier state carried forward from step
   $t-1$ --- the base weight remains frozen, every previously learned rank
   block is reinstated as frozen and active, and one new trainable block is
   appended (Section~\ref{sec:ch3-rankext}).
\item Update the classifier state for this step: rows for classes already
   introduced remain frozen, and rows for $C_t$ are unmasked and become
   trainable (Section~\ref{sec:ch3-classifier-construction}).
\item Train on $D_t$ with the classification objective and whichever
   stabilisation mechanisms are active for the selected variant
   (Section~\ref{sec:ch3-stabilization}) --- the distillation term against
   the frozen previous-step model and/or the FactorOrth term against the
   frozen accumulated prior blocks, each computed exactly as defined there.
\item At the end of the step, the adapter and classifier state --- now
   including the newly trained block and the newly trained classifier rows
   --- is carried forward as the retained state for step $t+1$; no raw
   example of $D_t$ is kept.
\end{enumerate}

Because the persistent state produced by this loop defines the model at
every point in the sequence, RankExt has no separate final-merge step
analogous to SimpleAvg's dense-averaging reconstruction; the model at the
end of step $T$ is the final model.

\textbf{Shared closing sequence.} For either family, once the final model is
assembled: the configured classifier rescaling is applied
(Section~\ref{sec:ch3-calibration}); open accuracy is computed from the resulting logits
(Section~\ref{sec:ch3-open-restricted}); restricted accuracy is computed from the identical logits
under the narrower per-group mask (Section~\ref{sec:ch3-open-restricted}); and the continual-learning
metrics defined next are computed from the accuracies this evaluation
produces.

\subsection{Backward Transfer and Forgetting Definitions}
\label{sec:ch3-bwt-forgetting}

Let $a_{i,j}$ denote open accuracy (Section~\ref{sec:ch3-open-restricted}) on class group $C_j$,
measured using the checkpoint produced after step $i$ has been trained.
$a_{i,j}$ is defined only for $j\le i$: a class group cannot be evaluated
before the step that introduces it has been trained. Its diagonal entries,
$a_{j,j}$, record each group's accuracy at the point it was first learned;
its final row, $a_{T,j}$ for $j=1,\dots,T$, records every group's accuracy
under the completed model. Backward transfer and forgetting are both computed from
this same matrix, but they summarise different properties of it and are
not interchangeable.

Backward transfer (BWT) compares the final model against each group's own
diagonal entry:

\begin{equation}
\label{eq:ch3-bwt}
\text{BWT} = \frac{1}{T-1}\sum_{j=1}^{T-1}\big(a_{T,j} - a_{j,j}\big),
\end{equation}

excluding the final group, for which no later comparison point exists
\parencite{lopezpaz2017gradient}. A positive value indicates that, on average,
later training left a group's accuracy higher than it was at its own
diagonal reference; a negative value indicates the opposite. The two
families do not supply this diagonal reference in the same way, and the
difference must be stated plainly rather than treated as a detail. For
RankExt, $a_{j,j}$ is a genuine checkpoint: the model defined by the
persistent state immediately after step $j$, evaluated on $C_j$. For SimpleAvg,
no single evolving model exists to checkpoint (Section~\ref{sec:ch3-simpleavg}), so
$a_{j,j}$ is substituted with step $j$'s own independently-trained
specialist's accuracy on $C_j$, at that specialist's best epoch, measured
before the eventual merge. This substitution is the closest available
equivalent and remains diagnostically useful, but it is not structurally
identical to RankExt's reference --- one is a real intermediate state of the
model whose later trajectory is being measured, the other is a separate
model that the final SimpleAvg model never itself passes through. BWT is
reported for both families under this understanding.

Forgetting instead compares each non-final group's \emph{best-ever} recorded
accuracy against its accuracy at the final checkpoint:

\begin{equation}
\label{eq:ch3-forgetting}
f_j = \max_{j\le s\le T} a_{s,j} \;-\; a_{T,j}, \qquad
\text{Forgetting} = \frac{1}{T-1}\sum_{j=1}^{T-1} f_j,
\end{equation}

following \cite{chaudhry2018riemannian}, where the maximum is taken over the
group's entire recorded trajectory, including the final checkpoint itself,
so that $f_j$ is non-negative by construction. Computing $f_j$ requires a genuine
accuracy trajectory for group $j$ across every later checkpoint, which only
RankExt's incrementally-evolving architecture provides; SimpleAvg's $T$
independently trained adapters are not sequential checkpoints of
one evolving model, so this trajectory does not exist for it. Forgetting is
therefore reported for RankExt only, and no substitute or approximate value
is computed or reported for SimpleAvg --- this is a scope limit of what the
architecture makes available, not an omission.

Backward transfer and forgetting are related but distinct, and one is not
simply the negation of the other. BWT compares only two points for each
group --- its diagonal reference and the final checkpoint --- so a group whose
accuracy rose after being learned, peaked, and then fell back exactly to
its starting level would show a BWT contribution of zero for that group,
since its diagonal and final values coincide. Forgetting instead compares the final checkpoint against
whichever intervening checkpoint scored that group highest, so the same
trajectory would register a strictly positive forgetting contribution,
because that peak is retained in the maximum even though it does not
appear at either endpoint BWT considers. The two metrics are reported,
where available, alongside the open and restricted accuracies of Section~\ref{sec:ch3-open-restricted},
not in place of them.
